
\section{Notebook 17: Higher-Order Transition Memory and Shuffle Baseline}

Notebook 17 tests whether normalized prime-gap residual states contain higher-order
transition memory beyond a first-order Markov operator.  Let $s_n$ denote a
quantile-discretized state of the normalized gap
\[
z_n = \frac{p_{n+1}-p_n}{\log p_n}.
\]
The first-order operator is
\[
P_{ij}=\Pr(s_{n+1}=j\mid s_n=i),
\]
while the empirical two-step operator is
\[
P^{(2)}_{ik}=\Pr(s_{n+2}=k\mid s_n=i).
\]
The Markov prediction is $\widehat{P}^{(2)}=P^2$, giving the residual
\[
\Delta^{(2)} = P^{(2)}-\widehat{P}^{(2)}.
\]

For this run, the mean two-step $L_2$ residual is
\[
\|\Delta^{(2)}\|_2 = 0.00793828,
\]
compared with a shuffle-baseline mean of
\[
0.00219103.
\]
The lag-one excess mutual information is
\[
I_1^{\rm real}-I_1^{\rm shuffle} =
0.0124681,
\]
and the lag-two excess mutual information is
\[
I_2^{\rm real}-I_2^{\rm shuffle} =
0.00211557.
\]
These diagnostics quantify ordered transition memory after controlling for the
one-point state distribution.
